% !PRINTED 58
\weberpage{58}

%<<CONT
tion would be finite in the whole $u$-plane and would therefore, by
§\,15, 8, have to be a constant.

If $\varphi(u)$ is a doubly periodic function with the periods
$\omega_1$, $\omega_2$, then the integral taken over the boundary of the period
 parallelogram is [because of (1)]

\[
\int d\log\varphi(u) = \int\limits_{u_0}^{u_0+\omega_1} [d\log\varphi(u) - d\log\varphi(u+\omega_2)]
\]
\[
{}-\int\limits_{u_0}^{u_0+\omega_2} [d\log\varphi(u) - d\log\varphi(u+\omega_1)] = 0\,.
\]

From this there follows, by §\,15, 6., the theorem:

2. \gesperrt{A doubly periodic function becomes zero in a
period parallelogram at just as many points as it becomes
infinite.} If $m$ is this number, then $m$ is called the order
of the doubly periodic function. By 1. there are no
doubly periodic functions of order zero.

If in the proof one replaces the function $\varphi(u)$ by $\varphi(u)-c$,
it follows that a doubly periodic function of order $m$
takes any arbitrary value $c$ at equally many points.

If we take the integral of the formula in 7., §\,15 over the boundary
 of the period parallelogram, there results:

\[
\int u\, d\log\varphi(u) = \int\limits_{u_0}^{u_0+\omega_1} [u\,d\log\varphi(u) - (u+\omega_2)\,d\log\varphi(u+\omega_2)]
\]
\[
{}-\int\limits_{u_0}^{u_0+\omega_2} [u\,d\log\varphi(u) - (u+\omega_1)\,d\log\varphi(u+\omega_1)]
\]
\[
{}=\omega_1\int\limits_{u_0}^{u_0+\omega_2} d\log\varphi(u) - \omega_2\int\limits_{u_0}^{u_0+\omega_1} d\log\varphi(u)\,.
\]

Now, because of the many-valuedness of the logarithm,

\[
\int\limits_{u_0}^{u_0+\omega_2} d\log\varphi(u) = \log\varphi(u+\omega_2) - \log\varphi(u) = 2\pi i\, m_1
\]
\[
\int\limits_{u_0}^{u_0+\omega_1} d\log\varphi(u) = \log\varphi(u+\omega_1) - \log\varphi(u) = -2\pi i\, m_2\,,
\]

where $m_1$ and $m_2$ are indeterminate integers.

If therefore the function $\varphi(u)$ becomes infinite at the points $\alpha_1$, $\alpha_2$, $\dots$, $\alpha_m$
and zero at the points $\beta_1$, $\beta_2$, $\dots$, $\beta_m$, then by
§\,15, 7.:

\[ \Sigma\alpha_i \equiv \Sigma\beta_i \ (\mathrm{mod}\ \omega_1,\ \omega_2)\,. \]
